\section[Struktur Halaman dengan Elemen Semantik]{Struktur Halaman dengan Elemen Semantik}

Bab 5 telah memperkenalkan elemen-elemen semantik HTML5 secara konseptual. Pada section ini, fokus diarahkan pada \textbf{praktik merancang} struktur halaman web menggunakan elemen-elemen tersebut. Struktur semantik yang baik bukan sekadar tentang tampilan visual, melainkan tentang memberikan \textbf{makna} pada setiap bagian halaman sehingga browser, mesin pencari, dan teknologi bantu dapat memahami peran masing-masing konten. Halaman yang terstruktur dengan baik lebih mudah dirawat, lebih aksesibel, dan mendapat peringkat SEO yang lebih baik \cite{webdevHtml}.

HTML5 menyediakan elemen semantik yang menggantikan penggunaan \code{<div>} generik untuk bagian-bagian halaman yang umum. Setiap elemen memiliki makna spesifik: \code{<header>} untuk kepala halaman atau section, \code{<nav>} untuk navigasi utama, \code{<main>} untuk konten utama (hanya satu per halaman), \code{<article>} untuk konten mandiri, \code{<section>} untuk pengelompokan tematik, \code{<aside>} untuk konten pendukung, dan \code{<footer>} untuk penutup halaman. Penggunaan elemen yang tepat membantu pembaca layar membuat peta navigasi otomatis bagi penggunanya \cite{mdnHtmlIntro}.

\begin{table}[htbp]
\centering
\begin{tabular}{|P{2.2cm}|P{4cm}|P{5.5cm}|}
\hline
\textbf{Elemen} & \textbf{Makna Semantik} & \textbf{Contoh Penggunaan} \\
\hline
\code{<header>} & Kepala halaman/section & Logo, judul situs, navigasi utama \\
\hline
\code{<nav>} & Blok navigasi & Menu utama, breadcrumb, sidebar link \\
\hline
\code{<main>} & Konten utama & Area isi utama (satu per halaman) \\
\hline
\code{<article>} & Konten mandiri & Blog post, berita, komentar \\
\hline
\code{<section>} & Grup tematik & Bab, tab panel, bagian konten \\
\hline
\code{<aside>} & Konten pendukung & Sidebar, iklan, info terkait \\
\hline
\code{<footer>} & Penutup halaman & Hak cipta, link kebijakan, kontak \\
\hline
\end{tabular}
\caption{Elemen semantik HTML5 dan konteks penggunaannya}
\end{table}

\begin{figure}[htbp]
\centering
\resizebox{0.85\textwidth}{!}{%
\begin{tikzpicture}[node distance=0.6cm, transform shape, every node/.style={font=\footnotesize}]
  % Body container
  \node (body) [class, minimum width=12cm, minimum height=0.8cm] {\code{<body>}};

  % Header
  \node (header) [class, below=of body, minimum width=11cm, fill=orange!15] {\code{<header>} --- Logo, Judul Situs};

  % Nav
  \node (nav) [class, below=of header, minimum width=11cm, fill=yellow!15] {\code{<nav>} --- Menu Navigasi Utama};

  % Main + Aside row
  \node (main) [class, below=of nav, minimum width=7.5cm, xshift=-1.75cm, minimum height=3cm, fill=green!10] {\code{<main>}};

  % Article inside main
  \node (article) [object, below=0.3cm of main.north, minimum width=6.5cm, fill=green!20] {\code{<article>} --- Konten Artikel};
  \node (section) [object, below=0.3cm of article, minimum width=6.5cm, fill=green!25] {\code{<section>} --- Bagian Konten};

  % Aside
  \node (aside) [class, right=0.5cm of main, minimum width=3cm, minimum height=3cm, fill=purple!10] {\code{<aside>}\\Sidebar};

  % Footer
  \node (footer) [class, below=0.6cm of main.south east, minimum width=11cm, xshift=-0.5cm, fill=gray!15] {\code{<footer>} --- Hak Cipta, Kontak};
\end{tikzpicture}%
}
\caption{Diagram struktur halaman web dengan elemen semantik HTML5}
\end{figure}

Hierarki heading (\code{<h1>} sampai \code{<h6>}) di dalam elemen semantik sangat penting untuk aksesibilitas dan SEO. Praktik terbaik: gunakan satu \code{<h1>} per halaman untuk judul utama, lalu \code{<h2>} untuk judul section, \code{<h3>} untuk sub-section, dan seterusnya. Jangan melewatkan level heading (misalnya langsung dari \code{<h2>} ke \code{<h4>}) karena akan membingungkan pembaca layar yang menggunakan heading untuk navigasi cepat. Setiap \code{<section>} dan \code{<article>} sebaiknya diawali dengan heading yang sesuai \cite{mdnAccessibility}.

\begin{webcode}[language=HTML, caption={Struktur halaman lengkap dengan elemen semantik}]
<!DOCTYPE html>
<html lang="id">
<head>
  <meta charset="UTF-8" />
  <title>Blog Pemrograman Web</title>
</head>
<body>
  <header>
    <h1>Blog Pemrograman Web</h1>
    <nav>
      <ul>
        <li><a href="/">Beranda</a></li>
        <li><a href="/artikel">Artikel</a></li>
        <li><a href="/tentang">Tentang</a></li>
      </ul>
    </nav>
  </header>

  <main>
    <article>
      <h2>Mengenal HTML5 Semantik</h2>
      <p>HTML5 memperkenalkan elemen semantik...</p>
      <section>
        <h3>Mengapa Semantik Penting?</h3>
        <p>Elemen semantik memberikan makna...</p>
      </section>
    </article>
    <aside>
      <h2>Artikel Terkait</h2>
      <ul>
        <li><a href="/css-dasar">CSS Dasar</a></li>
        <li><a href="/js-intro">Pengenalan JS</a></li>
      </ul>
    </aside>
  </main>

  <footer>
    <p>&copy; 2026 Blog Pemrograman Web</p>
  </footer>
</body>
</html>
\end{webcode}

\begin{catatan}
Jangan menggunakan elemen semantik hanya untuk mendapatkan styling default (misalnya \code{<nav>} agar terlihat seperti menu). Gunakan elemen berdasarkan \textbf{maknanya}, bukan tampilannya. Tampilan dapat selalu diubah dengan CSS. Sebaliknya, jangan menggunakan \code{<div>} untuk konten yang memiliki makna semantik yang jelas---gunakan elemen semantik yang tepat agar halaman Anda aksesibel dan mudah dirawat \cite{webdevHtml}.
\end{catatan}

